\section{The YAMSO Ledger Model}
\label{sec:model}

We describe the ledger interface at the level needed for the finite
construction in \cref{sec:basic-construction}. Our approach follows the
modelling style of \cite[Section~2.1]{DinsdaleYoungEtAl2020Afgjort}: we state
an intentionally strong and simple abstraction, isolate the properties used
by the protocol, and postpone its detailed realization. A later version will
express this interface as an ideal functionality
$\mathcal F_{\mathsf{Ledger}}$, in the style of the total-order-broadcast
functionality of \cite{DamgardEtAl2025AsyncYOSO}. There is no reason to fix
that UC bookkeeping before the API itself has stabilized.

The ledger is generic. It knows nothing about universal sampling, garbled
circuits, committees, input wires, or the meaning of protocol roles. Those
belong to the protocol using the ledger.

\subsection{Public state and initialization}

Fix a security parameter $\lambda$, a session identifier $\mathsf{sid}$, a
constant $p<1/3$, a generic setup distribution $\mathsf{Setup}$, a
key-generation algorithm $\mathsf{KeyGen}$, and a description of a deterministic
protocol-dependent algorithm $\mathsf{WhosTurn}$. The description of
$\mathsf{WhosTurn}$ is supplied to the ledger as a fixed initialization value;
it cannot be changed in response to the execution. The ledger samples
\[
  \mathsf{crs}\leftarrow\mathsf{Setup}(1^\lambda)
\]
and initializes the public append-only list as
$L=((\mathsf{Setup},\mathsf{crs}))$. The concrete protocol gives meaning to
$\mathsf{crs}$; the ledger does not interpret it.

We assume synchronous computation and an unrealistically strong broadcast
abstraction. After an entry is appended, every party sees exactly the same
list $L$ at the same time. An append is atomic, cannot be suppressed after
it occurs, and fixes the position of that entry. There are no delayed,
private, or party-local ledger views. This deliberately idealized assumption
removes consensus and network scheduling from the present question.

Apart from its own setup, registration, and timeout records, the ledger does
not enforce protocol syntax. In particular, an arbitrary message supplied by
a corrupt scheduled party is appended verbatim. The protocol reading $L$
decides which entries are well-formed and ignores malformed ones.

\subsection{Two registration paths}

The registry grows during the execution. There are two ways to add a key;
neither command takes a protocol role as input. The interface imposes no
a-priori bound on the registry, although any efficient execution creates only
polynomially many keys and every finite protocol instance uses a finite
prefix.

\paragraph{Computation-party registration.}
On a $\mathsf{Register}$ request from the adversary, the ledger samples
\[
  (pk_i,sk_i)\leftarrow\mathsf{KeyGen}(1^\lambda),
\]
where $i$ is the next computation-key index, creates an ITM whose PID is
$pk_i$, and appends $(\mathsf{Register},pk_i)$ to $L$. Independently it
samples $c_i\leftarrow\mathsf{Bernoulli}(p)$. If $c_i=1$, the new
computation party is corrupt and $sk_i$ is given to the adversary. Otherwise
it is honest, the ledger retains $sk_i$ as protected state, and the adversary
cannot corrupt this party later. Registration is indivisible: after
requesting it, the adversary cannot discard the public key because it
dislikes the corruption outcome.

\paragraph{Application-party registration.}
A fixed application party with ordinary UC identity $P$ may also input
$\mathsf{Register}$. The ledger samples
$(pk_P,sk_P)\leftarrow\mathsf{KeyGen}(1^\lambda)$, stores the association
$P\leftrightarrow(pk_P,sk_P)$, and appends
$(\mathsf{Register},P,pk_P)$ to $L$. It does not return $sk_P$ to an honest
$P$. There is no Bernoulli corruption experiment for such a party: the
adversary may corrupt $P$ adaptively in the ordinary UC sense. If $P$ is
already corrupt, or becomes corrupt later, the ledger releases $sk_P$ and
its retained record of $P$'s protected calls to the adversary. The target
ideal functionality supplies the simulator with any input or output leakage
prescribed upon that corruption.

Each PID registers at most once per session. Application parties are the
meaningful principals at the interface of the computation---for example, an
account owner supplying an input. Their identity need not and generally
cannot be hidden. Computation parties, by contrast, are protocol machinery.

A public protocol-level rule maps registration positions to roles. Thus the
seventh computation key may represent the eighth share-release role. This
mapping is not an argument to $\mathsf{Register}$ and is not interpreted by
the ledger.

\subsection{Turns and protected local evaluation}

At each synchronous step, the fixed algorithm $\mathsf{WhosTurn}(L)$ returns
either a registered public key that has not previously had a turn, or
$\bot$. It may schedule a computation key $pk_i$ or the registered key
$pk_P$ of an application party. A well-formed protocol chooses
$\mathsf{WhosTurn}$ so that it waits for enough registrations and schedules
each required role in the appropriate order. The ledger does not interpret
why a particular key is scheduled.

When a key $pk$ is scheduled, its owner has one synchronous opportunity to
speak. Write $L_{mathsf{pre}}$ for the list at the start of that turn. An
honest owner privately gives the ledger a finite polynomial-time algorithm
$A$. The ledger samples and retains any coins required by $A$, and evaluates
\[
  m\leftarrow A(sk,L_{\mathsf{pre}}),
\]
using the protected secret key associated with $pk$. If $m\neq\bot$, it
signs $(\mathsf{sid},L_{\mathsf{pre}},m)$ using the signing component of $sk$
and appends
\[
  (\mathsf{Message},pk,m,\sigma)
\]
to $L$. Neither $A$ nor the unrevealed part of $sk$ is output. The ledger
retains the algorithm and its coins as the protected call record rather than
modelling an erasure.

If the owner is corrupt, the adversary knows $sk$ and may supply any pair of
bit strings $(m,\sigma)$ during the turn. The ledger appends
$(\mathsf{Message},pk,m,\sigma)$ verbatim without checking a signature, a key
relation, or any other protocol-specific predicate. If an
honest $A$ returns $\bot$, or if no message is supplied by the end of the
turn, the ledger appends $(\mathsf{NoMsg},pk)$. This timeout is immediately
visible to everyone and lets $\mathsf{WhosTurn}$ advance past a silent role.
After either record is appended, $pk$ is spent and is never scheduled again.

For an honest nonempty message, protected evaluation, signing, and append are
one atomic operation. The adversary sees the result only after it is
irrevocably in $L$. The ledger does not give an adversarial scheduler an
opportunity to inspect the result and then suppress it.

We assume that $\mathsf{KeyGen}$ may provide a public relation
$\mathsf{KeyMatch}(pk,sk)$. A protocol that uses disclosure of $sk$ as its
message can then check that the disclosed key belongs to the scheduled
public key. If the chosen key system has no such relation, the ledger does
not manufacture one: the protocol must tolerate the missing check or choose
different key syntax.

\subsection{How protocols use the interface}

The following examples explain the intended use without adding commands to
the ledger.

\paragraph{Ordinary computation role.}
From $\mathsf{crs}$ and the public list, all parties may derive the same
protocol object, including a ciphertext $C_{pk}$ under a scheduled public
key. This object need not itself be an entry of $L$. The honest role supplies
an algorithm $A$ which derives $C_{pk}$, decrypts its local state with $sk$,
examines $L$, and returns $sk$ if the role is activated and $\bot$ otherwise.
If it returns $sk$, everyone can check $\mathsf{KeyMatch}$, decrypt
$C_{pk}$, and continue the public computation. Thus publishing the key is the
role's only message and models complete post-speech exposure. If $A$ returns
$\bot$, the secret key and encrypted role state remain hidden forever.

For this interpretation, every secret used by an ordinary role must be
recoverable from $sk$ and public protocol data. A protocol that gives the
role additional hidden auxiliary state cannot claim that publishing only
$sk$ models leakage of its complete state.

\paragraph{Strong input role.}
An application party $P$ embeds its private input in $A$ and causes the
ledger to append only the prescribed value $A(sk_P,L)$ and its signature.
For example, $A$ may decrypt two activation tokens and return the one selected
by an input bit. If $P$ remains honest, neither $sk_P$, the unused token, nor
the input-dependent algorithm is revealed. If $P$ is corrupted, the
adversary receives its ordinary UC state together with the key and protected
call history retained by the ledger; correspondingly the simulator receives
the input leakage specified by the target functionality. Corruption after
the turn does not permit a replacement message because $pk_P$ is already
spent.

The present finite construction has public output. A later private-output
variant may add a dual protected-read operation which evaluates a function of
$(sk_P,L)$ and returns only its result privately to $P$; no such operation is
needed here.

\subsection{Random corruption and finite committees}

The random coins apply only to computation-party registrations. Let $C$ be
any size-$k$ computation committee selected from registration indices without
using the hidden coins $(c_i)_i$, and let
$X=|\{i\in C:c_i=1\}|$. Then $X\sim\mathsf{Binomial}(k,p)$. At the concrete
choice $p=1/4$, Hoeffding's inequality gives
\[
  \Pr[X\geq k/3]\leq \exp(-k/72).
\]
If a finite execution uses $q$ such committees, a union bound gives
\[
  \Pr[\neg\mathsf{Good}]\leq q\exp(-k/72),
\]
where $\mathsf{Good}$ is the event that every committee has at most
$t=\lfloor(k-1)/3\rfloor$ corrupt members. Hence
$k=\Theta(\lambda+\log q)$ makes the bad event negligible for every
polynomial-size finite execution. Application parties do not enter this
calculation.

\subsection{Interpretation and limitations}

The preliminary experiment should not be read as saying that a real adversary
naturally corrupts machines by tossing independent coins. The intended
real-world picture has a large population of apparently interchangeable
machines and an adversary that may choose any fixed $p$ fraction of that
population to compromise. Fresh computation clients, keys, and role states
are placed on machines at random, without revealing the placement before a
machine is corrupted or speaks. A key's logical role may be public even though
the physical machine carrying its secret state is not linkable from the
ledger. Random placement therefore turns the adversary's fixed physical set
into an approximately random corrupt subset of each computation committee.
The independent Bernoulli coins in the present model collapse this placement
argument into the cleaner statement that each computation PID is corrupt
independently.

This is a fixed-population, non-mobile abstraction. Once a computation PID is
sampled honest it cannot later be corrupted. An adversary that can repeatedly
replace the corrupted quarter of the physical population is strictly stronger
and is not modelled here. Named application parties are different on purpose:
an account owner or other designated input provider is not plausibly hidden by
random placement, so it retains the ordinary adaptive corruption interface.
The random computation corruptions are an implementation assumption to be
simulated away, not a corruption rule for the target application.

The statement that the ledger stores $sk$ and evaluates $A(sk,L)$ is likewise
an ideal-interface device, not a proposal that a blockchain hold users' secret
keys or execute their private code. In the intended implementation, a machine
stores its own key and locally computes the one value it places on the public
board. Moving that local step behind the ledger interface prevents an honest
party's secret key from becoming an output visible to the UC environment. For
a computation role, publishing $sk$ is an explicit way to model that speaking
exposes all of its encrypted role state; for a named application party, only
the prescribed value is published unless the party is actually corrupted.

This abstraction is intentionally stronger and cleaner than such a system.
It assumes honest key generation, independent placement, no selective abort
of an unfavourable registration, instantaneous common knowledge of $L$,
unpreventable honest appends, exact synchronous timeouts, protected
evaluation under a secret key, and a protocol-dependent turn function fixed
at initialization. It does not model consensus, denial of registration, network
partitions, grinding, compromise correlations, or the mechanism that places
a computation key on a real machine. These assumptions must not be mistaken
for properties already obtained by a concrete blockchain.

Finally, the computation PIDs and their random corruption coins exist only in
the real or ledger-hybrid execution. The target ideal functionality has the
usual application interface: named input and output parties, their prescribed
adaptive corruptions, and the intended outputs. It contains no committees,
randomly corrupted computation population, garbled circuits, or
$\mathsf{WhosTurn}$. A security proof must simulate all of that protocol
machinery from the target functionality's leakage.

The final paper may keep only this interface and interpretation in the main
text. Once the API is stable, an appendix can formulate
$\mathcal F_{\mathsf{Ledger}}$ in UC notation and make precise ITM creation,
private delivery and retention of $A$, synchronous timeouts, and corruption
bookkeeping. Those details should express the API above rather than determine
it.
